\par\noindent\begin{minipage}{\linewidth}
% !TeX root = ../教材题库_全部条目_转写工作稿.tex
% Source: XJ-B1-C02-EX-P046-Y1243
% 第2章 / 2.2 从函数观点看一元二次方程 / 教材 p.46 / PDF 第53页
% 已对照原页转写题干及公式；待教师复核；未制作答案
% 原题及相邻原解截图（核对用）：../题目原图/XJ-B1-C02-EX-P046-Y1243-01.webp, ../题目原图/XJ-B1-C02-EX-P046-Y1243-02.webp
% 本文件只在转写完成后作为题干使用；OCR 条目尚不是可靠题干。
\begin{bankitem}{2.2 例题 例题1}{第2章 / 2.2 从函数观点看一元二次方程 / 教材 p.46 / PDF 第53页}{XJ-B1-C02-EX-P046-Y1243}{已对照原页转写题干及公式；待教师复核；未制作答案}
二次函数 $y=ax^2+bx+c$ 的图象如图 2.2-2 所示，顶点坐标为 $(2,2)$。根据图象解答下列问题：
\bankfig{C02-EX046}{0.38}
\subq{1}写出方程 $ax^2+bx+c=0$ 的两个根；
\subq{2}若方程 $ax^2+bx+c=k$ 有两个不相等的实数根，求 $k$ 的取值范围。
\end{bankitem}
\end{minipage}\par\medskip
